\documentclass[10pt,a4paper]{amsart}
\usepackage[margin=19mm]{geometry}
\usepackage{amsmath,amssymb,mathtools}
\usepackage[scr=rsfs,cal=euler]{mathalfa}
\usepackage{xfrac}
\usepackage[unicode,hidelinks,bookmarks=false]{hyperref}
\newcommand{\ie}{i.e.\ }
\newcommand{\cf}{cf.\ }
\newcommand{\resp}{respectively\ }
\newcommand{\Subsection}{\subsection}
\providecommand{\nobreakdash}{\nobreak\hbox{-}}
\setlength{\emergencystretch}{2em}
\multlinegap0pt
\usepackage{xcolor}
\usepackage[shortlabels]{enumitem}
\newtheorem{proposition}{Proposition}
\title{Answer-Distribution Trajectories: A Stochastic-Dynamics View of LLM Reasoning}
\author{Mar Gonzàlez I Català}
\date{Source: arXiv:2609.09030v1 (CC BY 4.0)}
\begin{document}
\maketitle

\noindent\emph{Source and licence.} Answer-Distribution Trajectories: A Stochastic-Dynamics View of LLM Reasoning. Version arXiv:2609.09030v1 (CC BY 4.0), distributed by arXiv under the Creative Commons Attribution 4.0 International licence, http://creativecommons.org/licenses/by/4.0/, as recorded in the arXiv licence metadata for this exact version and shown on the abstract page https://arxiv.org/abs/2609.09030v1. Full licence notice: \texttt{licenses/arxiv-2609.09030-CC-BY-4.0.txt}.\par\smallskip

\noindent\textit{Editorial scope.} Counted unit: the article's proposition prop:nonidentifiability (source0.tex lines 289--301). The article's complete original proof of it is delivered with it (source0.tex lines 713--763, one \texttt{proof} environment of 51 lines, which the article places away from the statement). No mathematics is written, repaired or re-derived: the statement and the proof are the article's own lines, in the article's own notation, and no retained line is substituted or re-flowed.  No further source-local context is needed: every cross-reference of every retained line resolves inside the two delivered spans.  Nothing was omitted from any retained span, so the package carries no omission record.  No retained line contains a citation, so the wrapper carries no bibliography and no bibliography entry.  No retained line contains a figure environment, an image-inclusion command or drawing code, so no image is delivered and the package declares no asset dependency.  Editorial formatting only, in addition: an AMS \texttt{amsart} preamble that restates the article's own declarations and macros used by the retained lines, namely the environments \texttt{proposition} on separate counters, together with the standard packages those lines need.  The retained environments print in this document as Proposition 1 in order of appearance, whereas the article prints the same items with the numbering of its own full document.  The complete native source of the article as distributed in the arXiv e-print for this exact version is bundled unchanged as \texttt{sources/209770/source0.tex}.
\begin{proposition}[Endpoint and entropy trajectory are non-identifying]
\label{prop:nonidentifiability}
Neither endpoint prediction nor entropy trajectory identifies the underlying answer-distribution dynamics. In particular:

\begin{enumerate}
    \item The endpoint map \(E\) is non-injective: there exist distinct answer-distribution trajectories \(\mathcal{T}\neq\mathcal{T}'\) such that
    $
    E(\mathcal{T})=E(\mathcal{T}').
    $
    Hence, endpoint prediction does not determine the sequence of predictive states that produced it.
    \item For \(|\mathcal{A}|\geq 2\), the entropy map \(\mathcal{H}\) is non-injective: there exist distinct answer-distribution trajectories \(\mathcal{T} \neq \mathcal{T}'\) such that \(\mathcal{H}(\mathcal{T}) = \mathcal{H}(\mathcal{T}')\). Hence, an entropy trajectory does not identify the underlying sequence of predictive states.
\end{enumerate}
\end{proposition}

\begin{proof}[Proof of Proposition~\ref{prop:nonidentifiability}]
Both claims can be witnessed on the binary answer space \(\mathcal{A}=\{a,b\}\).

For the first claim, consider
\[
\mathcal{T}_1:
(0.9,0.1)
\rightarrow
(0.9,0.1)
\rightarrow
(0.9,0.1),
\]
and
\[
\mathcal{T}_2:
(0.1,0.9)
\rightarrow
(0.4,0.6)
\rightarrow
(0.9,0.1).
\]
Both terminate with \(a\) as the dominant answer, and therefore
\[
E(\mathcal{T}_1)=E(\mathcal{T}_2)=a.
\]
Nevertheless, \(\mathcal{T}_1\) maintains the same dominant hypothesis throughout, whereas \(\mathcal{T}_2\) revises from \(b\) to \(a\). Hence the endpoint does not identify the preceding distribution dynamics.

For the second claim, take
\[
p=(0.9,0.1),
\qquad
q=(0.1,0.9),
\]
and consider the constant trajectories
\[
\mathcal{T}_1 : p \to p \to p,
\qquad
\mathcal{T}_2 : q \to q \to q.
\]
Since entropy is invariant under permutation of coordinates,
\[
H(p)=H(q),
\]
and therefore
\[
\mathcal{H}(\mathcal{T}_1)
=
\mathcal{H}(\mathcal{T}_2).
\]
However, \(\mathcal{T}_1 \neq \mathcal{T}_2\), since \(p\neq q\). Hence entropy profiles do not identify which hypotheses carry the predictive mass.
\end{proof}

\end{document}
